\documentclass[12pt,a4paper]{article}

\usepackage[a4paper,text={16.5cm,25.2cm},centering]{geometry}
\usepackage{lmodern}
\usepackage{amssymb,amsmath}
\usepackage{bm}
\usepackage{graphicx}
\usepackage{microtype}
\usepackage{hyperref}
\usepackage{minted}
\setlength{\parindent}{0pt}
\setlength{\parskip}{1.2ex}

\hypersetup
       {   pdfauthor = {  },
           pdftitle={  },
           colorlinks=TRUE,
           linkcolor=black,
           citecolor=blue,
           urlcolor=blue
       }






\begin{document}



\section{Matematično nihalo}
Tina Brdnik

\subsection{Opis naloge}
Naloga obravnava matematično nihalo, katerega odmik $kot(t)$ je podan z diferencialno enačbo:

\[
\ddot{\theta}(t) + \frac{g}{l}\sin(\theta(t)) = 0, \quad \theta(0) = \theta_0, \quad \dot{\theta}(0) = \omega_0.
\]
Enačbo drugega reda prevedemo na sistem prvega reda:

\[
\begin{cases}
\theta'(t) = \omega(t), \\
\omega'(t) = -\dfrac{g}{l}\sin(\theta(t)).
\end{cases}
\]
Za različne začetne pogoje primerjamo rešitev z nihanjem harmoničnega nihala, ki je podano z linearizirano enačbo:

\[
\ddot{\theta}(t) + \frac{g}{l}\theta(t) = 0.
\]
Pri harmoničnem nihalu je nihajni čas neodvisen od začetnega odmika, medtem ko pri matematičnem nihalu narašča z večjo začetno energijo.

\subsection{Opis rešitve}
Implementacija je razdeljena na več funkcij, združenih v modulu \texttt{Naloga3}:

\begin{itemize}
\item[1. ] \textbf{Metoda DOPRI5}  

\end{itemize}
Funkcija \texttt{dopri5} je lastna implementacija metode Dormand\ensuremath{\endash}Prince 5. reda (DOPRI5) z adaptivnim korakom.  Metoda uporablja oceno napake med približkom 5. in 4. reda ter na tej osnovi povečuje ali zmanjšuje dolžino koraka.  S tem zagotavlja visoko natančnost pri optimalnem številu korakov.

\begin{itemize}
\item[2. ] \textbf{Sistem matematičnega nihala}  

\end{itemize}
Enačbo drugega reda smo pretvorili v sistem prvega reda in zapisali funkcijo 'sistem\_nihalo(t,y)', ki vrača odvode $[\theta', \omega']$ za dano stanje.

\begin{itemize}
\item[3. ] \textbf{Energija sistema}  

\end{itemize}
Funkcija \texttt{energija\_nihalo(y)} izračuna mehansko energijo nihala kot:

\[
E = \tfrac{1}{2}\,\omega^2 + \frac{g}{l}\,(1 - \cos(\theta)).
\]
Energija se v točni rešitvi ohranja, zato je uporabna za preverjanje pravilnosti numerične metode.

\begin{itemize}
\item[4. ] \textbf{Reševanje nihala}  

\end{itemize}
Funkcija \texttt{resi\_nihalo(zacetni\_kot, zacetna\_kotna\_hitrost, cas\_konec)} uporabi metodo \texttt{dopri5} in funkcijo \texttt{sistem\_nihalo} za izračun rešitve na podanem časovnem intervalu.  Rezultat sta vektor časov in matrika, ki vsebuje vrednosti kota in kotne hitrosti v teh časih.

\subsection{Rezultati}
Rešili smo sistem enačb za matematično nihalo in rezultate primerjali s harmoničnim nihalom. 

Za manjše začetne kote se rešitvi skoraj popolnoma ujemata.  Če pa začetni kot povečamo, se pojavi razlika: matematično nihalo niha počasneje kot harmonično.

Na prvem grafu vidimo potek kota skozi čas za oba nihala.  Pri majhnih odmikov krivulji potekata skupaj, pri večjih pa se sčasoma vedno bolj razlikujeta.

Izračunali smo tudi nihajni čas matematičnega nihala za različne začetne kote. Vsakemu začetnemu kotu smo priredili začetno energijo in narisali graf odvisnosti nihajnega časa od energije.

Rezultati kažejo, da ima harmonično nihalo vedno enak nihajni čas, medtem ko pri matematičnem nihalu nihajni čas narašča z začetno energijo, kar je pričakovano.


\begin{minted}[texcomments = true, mathescape, fontsize=\small, xleftmargin=0.5em]{julia}
include("demo.jl")
plt = demo_nihalo()
display(plt)
\end{minted}
\includegraphics[width=\linewidth]{jl_XjkHJm/porocilo_0_1.pdf}


\end{document}
